\documentclass[../../main]{subfiles}

\begin{document}

\section{Tipos de Funções}
\label{sec:matematica:funcoes:tipos-de-funcoes}

\begin{defn}[Função Injetiva]
  \label{def:matematica:funcoes:tipos-de-funcoes:injetiva}

  Seja

  [
    f:A\to B.
  ]

  Diz-se que (f) é injetiva quando

  [
    (\forall x,y\in A)
    \bigl(
      f(x)=f(y)
      \Rightarrow
      x=y
    \bigr).
  ]

\end{defn}

\begin{obs}

  Uma função injetiva nunca identifica dois elementos distintos do
  domínio.

\end{obs}

\begin{ex}

  A função

  [
    f:\mathbb R\to\mathbb R,
    \qquad
    f(x)=2x+1
  ]

  é injetiva.

\end{ex}

\begin{prop}[Caracterização da Injetividade]
  \label{prop:matematica:funcoes:tipos-de-funcoes:caracterizacao-injetiva}

  Seja

  [
    f:A\to B.
  ]

  Então (f) é injetiva se, e somente se,

  [
    (\forall x,y\in A)
    \bigl(
      x\neq y
      \Rightarrow
      f(x)\neq f(y)
    \bigr).
  ]

\end{prop}

\begin{proof}

  A afirmação é a contrapositiva da definição de injetividade.

\end{proof}

\begin{defn}[Função Sobrejetiva]
  \label{def:matematica:funcoes:tipos-de-funcoes:sobrejetiva}

  Seja

  [
    f:A\to B.
  ]

  Diz-se que (f) é sobrejetiva quando

  [
    (\forall y\in B)
    (\exists x\in A)
    \bigl(
      f(x)=y
    \bigr).
  ]

\end{defn}

\begin{obs}

  Uma função é sobrejetiva quando sua imagem coincide com seu
  contradomínio.

\end{obs}

\begin{prop}[Caracterização da Sobrejetividade]
  \label{prop:matematica:funcoes:tipos-de-funcoes:caracterizacao-sobrejetiva}

  Seja

  [
    f:A\to B.
  ]

  Então

  [
    f
    \text{ é sobrejetiva}
    \iff
    \operatorname{Im}(f)=B.
  ]

\end{prop}

\begin{proof}

  Segue diretamente das definições de imagem e sobrejetividade.

\end{proof}

\begin{ex}

  A função

  [
    f:\mathbb R\to\mathbb R,
    \qquad
    f(x)=x^3
  ]

  é sobrejetiva.

\end{ex}

\begin{defn}[Função Bijetiva]
  \label{def:matematica:funcoes:tipos-de-funcoes:bijetiva}

  Seja

  [
    f:A\to B.
  ]

  Diz-se que (f) é bijetiva quando é simultaneamente injetiva e
  sobrejetiva.

\end{defn}

\begin{ex}

  A função

  [
    f:\mathbb R\to\mathbb R,
    \qquad
    f(x)=x+1
  ]

  é bijetiva.

\end{ex}

\begin{prop}
  \label{prop:matematica:funcoes:tipos-de-funcoes:bijetiva-unica}

  Seja

  [
    f:A\to B.
  ]

  Se (f) é bijetiva, então para cada
  (y\in B) existe um único elemento
  (x\in A) tal que

  [
    f(x)=y.
  ]

\end{prop}

\begin{proof}

  A existência decorre da sobrejetividade.

  A unicidade decorre da injetividade.

\end{proof}

\subsection{Composição}

\begin{prop}[Composição de Injetivas]
  \label{prop:matematica:funcoes:tipos-de-funcoes:composicao-injetiva}

  Se

  [
    f:A\to B
  ]

  e

  [
    g:B\to C
  ]

  são injetivas, então

  [
    g\circ f
  ]

  é injetiva.

\end{prop}

\begin{proof}

  Suponha

  [
    (g\circ f)(x)
    =
    (g\circ f)(y).
  ]

  Então

  [
    g(f(x))
    =
    g(f(y)).
  ]

  Como (g) é injetiva,

  [
    f(x)=f(y).
  ]

  Como (f) é injetiva,

  [
    x=y.
  ]

\end{proof}

\begin{prop}[Composição de Sobrejetivas]
  \label{prop:matematica:funcoes:tipos-de-funcoes:composicao-sobrejetiva}

  Se

  [
    f:A\to B
  ]

  e

  [
    g:B\to C
  ]

  são sobrejetivas, então

  [
    g\circ f
  ]

  é sobrejetiva.

\end{prop}

\begin{proof}

  Seja (c\in C).

  Como (g) é sobrejetiva, existe
  (b\in B) tal que

  [
    g(b)=c.
  ]

  Como (f) é sobrejetiva, existe
  (a\in A) tal que

  [
    f(a)=b.
  ]

  Logo

  [
    (g\circ f)(a)
    =
    c.
  ]

\end{proof}

\begin{prop}[Composição de Bijetivas]
  \label{prop:matematica:funcoes:tipos-de-funcoes:composicao-bijetiva}

  A composição de funções bijetivas é bijetiva.

\end{prop}

\begin{proof}

  Segue das duas proposições anteriores.

\end{proof}

\subsection{Endofunções}

\begin{defn}[Endofunção]
  \label{def:matematica:funcoes:tipos-de-funcoes:endofuncao}

  Uma endofunção é uma função

  [
    f:A\to A.
  ]

\end{defn}

\begin{ex}

  A função identidade

  [
    \operatorname{id}_A:A\to A
  ]

  é uma endofunção.

\end{ex}

\begin{defn}[Permutação]
  \label{def:matematica:funcoes:tipos-de-funcoes:permutacao}

  Uma permutação de um conjunto (A) é uma bijeção

  [
    f:A\to A.
  ]

\end{defn}

\begin{obs}

  O conjunto de todas as permutações de (A) será importante
  posteriormente em Álgebra Abstrata, onde forma um grupo sob a
  operação de composição.

\end{obs}

\end{document}
